\documentclass[a4paper,11pt]{article}
\usepackage[T2A]{fontenc}
\usepackage[utf8]{inputenc}
\usepackage[russian]{babel}
\usepackage{amsmath,amssymb,mathtools}
\usepackage{geometry}
\geometry{left=23mm,right=23mm,top=24mm,bottom=23mm}
\usepackage{microtype}
\usepackage{icomma}
\usepackage{booktabs,tabularx,array}
\usepackage{enumitem}
\usepackage{tikz}
\usepackage{tkz-euclide}
\usetikzlibrary{arrows.meta,positioning,shapes.geometric,calc}
\usepackage{xcolor}
\usepackage{fancyhdr}
\usepackage{hyperref}
\definecolor{accent}{HTML}{087D81}
\definecolor{darkaccent}{HTML}{075D60}
\definecolor{textgray}{HTML}{424242}
\hypersetup{hidelinks,pdftitle={Чек-лист. Стереометрия: куб, параллелепипед и призма},pdfauthor={Лёвин Артём Александрович}}
\setlength{\parindent}{0pt}
\setlength{\parskip}{6.5pt}
\linespread{1.10}
\setlist[itemize]{leftmargin=5.5mm,itemsep=2.5pt,topsep=2.5pt}
\setlist[enumerate]{leftmargin=7mm,itemsep=5pt,topsep=4pt}
\setlength{\headheight}{22pt}
\pagestyle{fancy}
\fancyhf{}
\fancyhead[L]{\small\color{textgray}Стереометрия: объёмы и площади}
\fancyhead[R]{\small\color{textgray}Николь Саркисьянц}
\renewcommand{\headrulewidth}{0.35pt}
\renewcommand{\headrule}{\hbox to\headwidth{\color{accent}\leaders\hrule height \headrulewidth\hfill}}
\fancyfoot[C]{\small\thepage}
\newcommand{\key}[1]{{\color{darkaccent}\bfseries #1}}
\newcommand{\example}[1]{\par\vspace{2pt}{\color{darkaccent}\bfseries #1}\par\vspace{-1pt}}
\newcommand{\sect}[1]{\section{#1}\vspace{-4pt}}
\renewcommand{\thesection}{\arabic{section}}
\begin{document}
\thispagestyle{empty}
\begin{center}
\vspace*{9mm}
{\color{darkaccent}\Large МАТЕМАТИКА}\par
\vspace{28mm}
{\huge Чек-лист}\par
\vspace{14mm}
{\LARGE\color{accent} Стереометрия}\par
\vspace{5mm}
{\large\color{darkaccent} Куб, прямоугольный параллелепипед,}\par
{\large\color{darkaccent} прямая треугольная призма}\par
\vspace{20mm}
{\large Подготовка к ЕГЭ по профильной математике}\par
\vfill
{\large Николь Саркисьянц}\par
\vspace{13mm}
{\large 9 ноября 2026 г.}\par
\vspace{20mm}
{Лёвин Артём Александрович}\\[2pt]
{эксклюзивно для Николь Саркисьянц}\par
\vspace{19mm}
\begin{tikzpicture}[x=1cm,y=1cm]
\draw[accent,line width=.85pt] (0,0) .. controls (1.5,.65) and (3,-.6) .. (4.8,0)
 .. controls (6.4,.65) and (7.2,-.6) .. (9,0);
\fill[darkaccent] (4.8,0) circle (1.5pt);
\end{tikzpicture}
\vspace{2mm}
\end{center}
\clearpage

\sect{Объём, площадь поверхности и высота}

\key{Объём} показывает, сколько места занимает тело в пространстве; измеряется в кубических единицах. \key{Площадь поверхности}~--- сумма площадей всех его граней; измеряется в квадратных единицах.

Для \emph{прямой призмы}, включая прямоугольный параллелепипед, работают три основные формулы:
\[
\boxed{V=S_0h},\qquad
\boxed{S_{\text{бок}}=P_0h},\qquad
\boxed{S_{\text{полн}}=2S_0+P_0h}.
\]
Здесь $S_0$~--- площадь одного основания, $P_0$~--- периметр основания, $h$~--- высота (длина бокового ребра у прямой призмы).

\key{Почему так считается боковая поверхность?} Если развернуть боковые грани прямой призмы в одну полосу, получится прямоугольник шириной $P_0$ и высотой $h$.

\begin{center}
\begin{tikzpicture}[scale=.83,line width=.85pt]
% Cuboid with projected orthogonal edges
\tkzDefPoint(0,0){A}
\tkzDefPoint(3.1,0){B}
\tkzDefPoint(4.25,1.05){C}
\tkzDefPoint(1.15,1.05){D}
\tkzDefPoint(0,2.4){At}
\tkzDefPoint(3.1,2.4){Bt}
\tkzDefPoint(4.25,3.45){Ct}
\tkzDefPoint(1.15,3.45){Dt}
\tkzDrawSegments(A,B B,C C,Ct Ct,Bt Bt,At At,A At,Dt Dt,Ct Bt,Ct)
\tkzDrawSegments[dashed,color=gray](A,D D,C D,Dt)
\node[below] at (1.55,0) {$a$};
\node[right] at (3.78,.47) {$b$};
\node[left] at (0,1.2) {$h=c$};
\node[text=darkaccent,align=center] at (7.3,2.65) {Прямоугольный\\параллелепипед};
\node[align=left] at (7.2,1.4) {$S_0=ab$\\ $P_0=2(a+b)$\\ $V=abh$};
\end{tikzpicture}
\end{center}

\key{Памятка по основанию.}
\begin{center}
\renewcommand{\arraystretch}{1.3}
\begin{tabularx}{\linewidth}{@{}p{53mm}XX@{}}
\toprule
Основание & Площадь $S_0$ & Периметр $P_0$\\
\midrule
Прямоугольник $a\times b$ & $ab$ & $2(a+b)$\\
Прямоугольный треугольник с катетами $a,b$ & $\dfrac{ab}{2}$ & $a+b+\sqrt{a^2+b^2}$\\
Равносторонний треугольник со стороной $a$ & $\dfrac{a^2\sqrt3}{4}$ & $3a$\\
\bottomrule
\end{tabularx}
\end{center}

\key{Самопроверка.} В формуле площади \emph{полной} поверхности присутствуют два основания: $2S_0$. Если дано сечение, перпендикулярное ребру прямоугольного параллелепипеда, его площадь равна произведению двух других рёбер.

\clearpage
\sect{Прямоугольный параллелепипед}

Пусть три взаимно перпендикулярных ребра из одной вершины равны $a,b,c>0$.
\[
\boxed{V=abc},\quad
\boxed{S_{\text{полн}}=2(ab+ac+bc)}.
\]
Для пространственной диагонали $d$ и диагонали прямоугольной грани $d_{ab}$:
\[
\boxed{d^2=a^2+b^2+c^2},\qquad
\boxed{d_{ab}^2=a^2+b^2}.
\]

\key{Две теоремы Пифагора.} Сначала у прямоугольной грани найдите $d_{ab}^2=a^2+b^2$. Затем рассмотрите прямоугольный треугольник с катетами $d_{ab}$ и $c$: $d^2=d_{ab}^2+c^2$.

\example{Пример 1. Площадь сечения по объёму}
Объём параллелепипеда равен $24$, одно ребро имеет длину $3$. Найдите площадь грани, перпендикулярной этому ребру.
\[
V=S_0h,\qquad S_0=\frac{V}{h}=\frac{24}{3}=8.
\]
\key{Ответ:} $8$ квадратных единиц.

\example{Пример 2. Неизвестное ребро по диагонали}
Два ребра равны $2$ и $4$, пространственная диагональ равна $6$. Найдите третье ребро и объём.
\[
6^2=2^2+4^2+c^2\ \Longrightarrow\ c^2=36-4-16=16.
\]
Так как длина ребра положительна, $c=4$. Следовательно,
\[
V=2\cdot4\cdot4=\boxed{32}.
\]

\key{Типичная ошибка.} Пространственная диагональ учитывает \emph{три} взаимно перпендикулярных измерения. Выражение $a^2+b^2$ подходит лишь для диагонали одной грани.

\clearpage
\sect{Куб: частный случай параллелепипеда}

У куба все рёбра равны $a>0$, а каждая грань~--- квадрат.
\[
\boxed{V=a^3},\qquad\boxed{S_{\text{полн}}=6a^2},\qquad
\boxed{d=a\sqrt3},\qquad\boxed{d_{\text{грани}}=a\sqrt2}.
\]
\example{Пример 3. Ребро куба увеличили на единицу}
Если ребро увеличить на $1$, объём увеличится на $19$. Найдите первоначальное ребро $a$.
\begin{align*}
(a+1)^3-a^3&=19,\\
a^3+3a^2+3a+1-a^3&=19,\\
3a^2+3a-18&=0,\\
a^2+a-6&=0,\\
(a+3)(a-2)&=0.
\end{align*}
Из двух формальных корней $-3$ и $2$ длине ребра соответствует только $a=2$.\par
\key{Проверка:} $3^3-2^3=27-8=19$.

\example{Пример 4. Как изменяется объём при масштабировании}
Если ребро куба увеличить втрое, то
\[
\frac{V_{\text{нов}}}{V_{\text{стар}}}
=\frac{(3a)^3}{a^3}=\boxed{27}.
\]
При одновременном увеличении \emph{всех линейных размеров} в $m>0$ раз:
\[
\boxed{\frac{S_{\text{нов}}}{S_{\text{стар}}}=m^2},\qquad
\boxed{\frac{V_{\text{нов}}}{V_{\text{стар}}}=m^3}.
\]
Это относится и к прямоугольному параллелепипеду, и к призме при равномерном масштабировании.

\key{Не путайте.} Увеличить ребро \emph{на} $1$ и увеличить его \emph{в} $3$ раза~--- разные действия. В первом случае сравниваются $(a+1)^3$ и $a^3$, во втором~--- $(3a)^3$ и $a^3$.

\clearpage
\sect{Прямая треугольная призма}

В основаниях треугольной призмы лежат два равных параллельных треугольника. У \emph{прямой} призмы боковые рёбра перпендикулярны основанию, поэтому их длина равна высоте $h$.
\[
\boxed{V=S_0h},\qquad
\boxed{S_{\text{полн}}=2S_0+P_0h}.
\]
\begin{center}
\begin{tikzpicture}[scale=.88,line width=.85pt]
\tkzDefPoint(0,0){A}
\tkzDefPoint(3.4,0){B}
\tkzDefPoint(2.3,1.5){C}
\tkzDefPoint(0,2.65){At}
\tkzDefPoint(3.4,2.65){Bt}
\tkzDefPoint(2.3,4.15){Ct}
\tkzDrawSegments(A,B B,C A,At B,Bt At,Bt Bt,Ct Ct,At)
\tkzDrawSegments[dashed,color=gray](A,C C,Ct)
\node[left] at (0,1.3) {$h$};
\node[below] at (1.7,0) {$a$};
\node[right] at (3.05,.75) {$b$};
\node[darkaccent,align=left] at (6.8,2.4) {$S_0$~--- площадь\\ треугольника\\ $P_0$~--- его периметр};
\end{tikzpicture}
\end{center}

\example{Пример 5. Все рёбра правильной треугольной призмы равны $\sqrt3$}
Основание~--- равносторонний треугольник со стороной $a=\sqrt3$. Боковое ребро $h=\sqrt3$.
\[
S_0=\frac{a^2\sqrt3}{4}=\frac{3\sqrt3}{4},\qquad
V=S_0h=\frac{3\sqrt3}{4}\cdot\sqrt3=\boxed{\frac94}.
\]
\key{Важно:} формулу $S_0=a^2\sqrt3/4$ применяйте только к \emph{равностороннему} треугольнику.

\example{Пример 6. Высота по объёму}
В основании лежит прямоугольный треугольник с катетами $3$ и $5$. Объём прямой призмы равен $30$.
\[
S_0=\frac{3\cdot5}{2}=\frac{15}{2},\qquad
h=\frac{30}{15/2}=\boxed4.
\]

\clearpage
\sect{Полная поверхность треугольной призмы}

\example{Пример 7. Основание с катетами $6$ и $8$}
Высота прямой призмы равна $10$. Найдите площадь полной поверхности.

\textbf{1.} Площадь основания:
\[
S_0=\frac{6\cdot8}{2}=24.
\]
\textbf{2.} Гипотенуза основания по теореме Пифагора:
\[
c=\sqrt{6^2+8^2}=\sqrt{100}=10.
\]
\textbf{3.} Периметр основания: $P_0=6+8+10=24$.

\textbf{4.} Площадь боковой и полной поверхности:
\[
S_{\text{бок}}=24\cdot10=240,\qquad
S_{\text{полн}}=240+2\cdot24=\boxed{288}.
\]

\subsection*{Одна формула, разные искомые величины}
Если известно $V$, то по формуле $V=S_0h$ можно найти любую неизвестную величину:
\[
\boxed{h=\frac{V}{S_0}},\qquad\boxed{S_0=\frac{V}{h}}.
\]
Если известна полная поверхность прямой призмы, а основание уже найдено, то
\[
\boxed{h=\frac{S_{\text{полн}}-2S_0}{P_0}},\qquad P_0>0.
\]

\subsection*{Точки контроля}
\begin{itemize}
\item Правильно определите вид основания: прямоугольник, прямоугольный или равносторонний треугольник.
\item При вычислении боковой поверхности используйте \emph{периметр} основания, при вычислении объёма~--- его \emph{площадь}.
\item Для полной поверхности добавьте площади \emph{двух} оснований.
\item Высота прямой призмы равна боковому ребру; у другой призмы это утверждение может быть неверным.
\item Не допускайте отрицательных длин и сверяйте размерность ответа.
\end{itemize}

\clearpage
\section*{Тренировочный блок~--- Путь А}
\textcolor{darkaccent}{\bfseries 10 задач.} Выполняйте задания последовательно: от базовых вычислений к комбинированным задачам. Ответы расположены на следующей странице.

\begin{enumerate}
\item Ребро куба равно $4$. Найдите его объём и площадь полной поверхности.
\item Рёбра прямоугольного параллелепипеда равны $2$, $3$ и $6$. Найдите объём, площадь полной поверхности и пространственную диагональ.
\item Объём прямоугольного параллелепипеда равен $168$, одно из его рёбер равно $7$. Найдите площадь грани, перпендикулярной этому ребру.
\item Ребро куба равно $5$. Найдите диагональ грани и пространственную диагональ. Во сколько раз увеличатся площадь полной поверхности и объём, если каждое ребро удвоить?
\item Если ребро куба увеличить на $1$, объём увеличится на $37$. Найдите первоначальную длину ребра.
\item Два ребра прямоугольного параллелепипеда равны $3$ и $4$, пространственная диагональ равна $13$. Найдите третье ребро, объём и площадь полной поверхности.
\item В основании прямой треугольной призмы лежит прямоугольный треугольник с катетами $3$ и $4$, высота призмы равна $10$. Найдите объём и площадь полной поверхности.
\item Все рёбра правильной треугольной призмы равны $2\sqrt3$. Найдите объём и площадь полной поверхности.
\item В основании прямой треугольной призмы лежит прямоугольный треугольник с катетами $5$ и $12$. Объём призмы равен $180$. Найдите её высоту и площадь полной поверхности.
\item Основание прямой треугольной призмы~--- прямоугольный треугольник с катетами $6$ и $8$. Площадь полной поверхности призмы равна $288$. Найдите высоту и объём.
\end{enumerate}

\clearpage
\section*{Ответы}
\vspace{5mm}
\begin{center}
\renewcommand{\arraystretch}{1.48}
\begin{tabularx}{\linewidth}{@{}cX@{}}
\toprule
№ & Ответ\\
\midrule
1 & $V=64$, $S_{\text{полн}}=96$.\\
2 & $V=36$, $S_{\text{полн}}=72$, $d=7$.\\
3 & $24$.\\
4 & $d_{\text{грани}}=5\sqrt2$, $d=5\sqrt3$; площадь~--- в $4$ раза, объём~--- в $8$ раз.\\
5 & $3$.\\
6 & Третье ребро $12$, $V=144$, $S_{\text{полн}}=192$.\\
7 & $V=60$, $S_{\text{полн}}=132$.\\
8 & $V=18$, $S_{\text{полн}}=36+6\sqrt3$.\\
9 & $h=6$, $S_{\text{полн}}=240$.\\
10 & $h=10$, $V=240$.\\
\bottomrule
\end{tabularx}
\end{center}
\vspace{15mm}
\section*{Итоговая памятка}
\begin{tabularx}{\linewidth}{@{}p{43mm}X@{}}
\toprule
Что дано & Полезное первое действие\\
\midrule
Куб, ребро $a$ & Используйте $V=a^3$, $S_{\text{полн}}=6a^2$.\\
Параллелепипед, три ребра & $V=abc$; для диагонали $d^2=a^2+b^2+c^2$.\\
Куб, изменение ребра & Сравните $a^3$ с $(a+t)^3$ или с $(ma)^3$.\\
Прямая треугольная призма & Найдите $S_0$ и $P_0$ её треугольного основания.\\
Площадь полной поверхности & $S_{\text{полн}}=P_0h+2S_0$.\\
\bottomrule
\end{tabularx}
\vfill
\begin{center}
\small\color{textgray}Лёвин Артём Александрович эксклюзивно для Николь Саркисьянц
\end{center}
\end{document}
